\documentclass[slide05.tex]{subfiles}
\begin{document}
\begin{frame}[plain]
    \begin{center}
        \vspace{48pt}
        {\huge\bf 赤外線（せきがいせん）}
    \end{center}
\end{frame}

\begin{frame}[fragile]
    \frametitle{赤外線(IR)とは}
    \vspace{20pt}
    \begin{center}
        \includesvg[width=\linewidth]{../../asset/image/electromagnetic_wave.svg}
    \end{center}
    \textpageref{26}
\end{frame}

\begin{frame}
    \frametitle{赤外線で通信する}
    \begin{itemize}
        \item ついていたらA, 消えていたらB
        \item 例
        \begin{itemize}
            \item 点いている: ロボットが前に進む
            \item 消えている: ロボットが止まる
        \end{itemize}
    \end{itemize}
    % \note{画像を入れる}
    \textpageref{26}
\end{frame}

\begin{frame}
   \frametitle{2通りでは足りない} 
   \begin{itemize}
        \item 2通りでは伝えられることが少ない
        \begin{itemize}
            \item 止まる
            \item 前に進む
        \end{itemize}
        \item もっといろいろなことを伝えたい
        \begin{itemize}
            \item 左に進む
            \item 右に進む
            \item 飛ぶ
            \item \dots
        \end{itemize}
   \end{itemize}
   \textpageref{26}
\end{frame}

\begin{frame}
    \frametitle{リモコンの原理} 
    \vspace{15pt}
    \begin{itemize}
        \item 1秒点いて2秒消えてまた2秒点いたら歩く
        \item 2秒点いて3秒消えたらジャンプ
    \end{itemize}
    \begin{center}
        \includesvg[width=0.8\linewidth]{../../asset/image/text05-img037.svg}
    \end{center}
    \textpageref{27}
\end{frame}

\begin{frame}
    \frametitle{赤外線をつかった家電}
    \begin{figure}[H]
    \begin{minipage}[t]{0.3\columnwidth}
        \centering
    \includesvg[scale=0.1]{../../asset/image/tv.svg}
        \caption{テレビ}
    \end{minipage}
    %\hspace{0.01\columnwidth} % ここで隙間作成
    \begin{minipage}[t]{0.3\columnwidth}
        \centering
        \includesvg[scale=0.1]{../../asset/image/aircondition.svg}
        \caption{エアコン}
    \end{minipage}
    \begin{minipage}[t]{0.3\columnwidth}
        \centering
        \includesvg[scale=0.1]{../../asset/image/fan.svg}
        \caption{扇風機}
    \end{minipage}
    \end{figure}
    \textpageref{27}
\end{frame}

\begin{frame}
    \frametitle{Raspberry Piをリモコンにする}
    \begin{itemize}
        \item リモコンを1つにまとめられる
        \vspace{10pt}
        \item プログラムで\ruby{制御}{せい|ぎょ}できる
        \begin{itemize}
            \item 外出先からエアコンをつける
            \item 音声でテレビのチャンネルを変える
        \end{itemize}
    \end{itemize}
    \textpageref{28}
\end{frame}

\begin{frame}
    \frametitle{Raspberry Piリモコンのしくみ}
    \vspace{12pt}
    \begin{center}
        \includesvg[width=0.7\linewidth]{../../asset/image/IRproc.svg}
    \end{center}
    \begin{textblock*}{0.4\linewidth}(90pt,180pt)
        \begin{tikzpicture}
            \node[rectangle,fill=cyan!10,text width=4cm, text centered,rounded corners]{赤外線の\ruby{点滅}{てん|めつ}で\\情報を伝える};
        \end{tikzpicture}
    \end{textblock*}
    \textpageref{28}
\end{frame}

\begin{frame}
    \frametitle{ラズパイをリモコンにする手順} 
    リモコンの動作をマネすることが必要
    \begin{enumerate}
        \item リモコンの信号を受信してファイルに記録する
        \item ファイルに記録された信号を赤外線LEDから出力する
    \end{enumerate}
    \textpageref{29}
\end{frame}

\begin{frame}
    \frametitle{チーム分け}
    班ごとに
    \begin{itemize}
        \item ロボットを使うチーム
        \item \ruby{扇風機}{せん|ぷう|き}を使うチーム
    \end{itemize}
    で分かれて、ロボットか扇風機を準備しよう
\end{frame}

\begin{frame}
    \frametitle{ラズパイをリモコンにする -1} 
    \begin{enumerate}
        \item cd \textasciitilde/05
        \item sudo service lircd stop
        \item mode2 -d /dev/lirc1 \textbar tee onoff.txt
        \begin{itemize}
            \item 受信機に向かってリモコンのボタンを押そう
            \item onoff.txtに受信した信号が記録される
            \item Ctrl+cキーで終了！
        \end{itemize}
        \item convert\_pattern onoff.txt \textgreater \space onoff.pattern
    \end{enumerate}
    \textpageref{29}
\end{frame}

\begin{frame}
    \frametitle{ラズパイをリモコンにする -2} 
    onnoff.patternを使って設定ファイルを作る
    \begin{enumerate}
        \item cp template.lircd.conf 05.lircd.conf
        \item mousepad 05.lircd.conf
        \item mousepad onoff.pattern
    \end{enumerate}
    \textpageref{29}
\end{frame}

\begin{frame}[fragile]
    \frametitle{ラズパイをリモコンにする -3}
    赤字の部分を書き換えよう
    \begin{lstlisting}
    begin remote
        name <#red#fan#>
        flags RAW_CODES
        eps 30
        aeps 100

        gap 200000
        toggle_bit_mask 0x0
    
        begin raw_codes
            name <#red#onoff#>
            <#red#1207 588 447 838 447 823 475 812 1218 ...#>
        end raw_codes
    end remote
    \end{lstlisting}
    \textpageref{29}
\end{frame}

\begin{frame}
    \frametitle{ラズパイをリモコンにする -4}
    \begin{itemize}
        \item sudo cp 05.lircd.conf /etc/lirc/lircd.conf.d/
        \begin{itemize}
            \item 設定ファイルをコピー
        \end{itemize}
        \item sudo service lircd restart
        \begin{itemize}
            \item 設定ファイルが再読み込みされる
        \end{itemize}
        \item irsend SEND\_ONCE \color{red}fan onoff\color{black}
    \end{itemize}
    \textpageref{29-30}
\end{frame}

\begin{frame}
    \begin{exampleblock}{問題をといてみよう}
    \begin{itemize}
        \item 教科書30ページ 例題5-12(1問)
        \begin{itemize}
            \item 手順を読みながら、赤外線を入力・出力してみよう
        \end{itemize}
    \end{itemize}
    \end{exampleblock}
\end{frame}

\begin{frame}
    \frametitle{HSPで赤外線を出力}
    \begin{itemize}
        \item exec命令を使ってirsendコマンドを呼び出す
        \begin{itemize}
            \item exec "irsend SEND\_ONCE fan onoff" 
        \end{itemize}
    \end{itemize}
    \textpageref{31}
\end{frame}

\begin{frame}[fragile]
    \frametitle{ボタンを押して扇風機を回す}
    \begin{lstlisting}[title=\textasciitilde/05/ir.hsp]
    #include "hsp3dish.as"
    BUTTON_PIN = 24
    prev = 0
    status=0

    *main
        redraw 0 
        font "",20
        pos 0,0
    mes "ボタンPUSHで赤外線照射!"
        redraw 1
        wait 1
        goto *edge

    *edge
        current = gpioin(BUTTON_PIN)
        if (prev=0) & (current=1) {	 <#blue#;ボタンが押されたら#>
            exec "irsend SEND_ONCE fan onoff"<#blue#;irsendコマンドを実行する#>
        }
        prev = current
        wait 1
        goto *edge
    \end{lstlisting}
    \textpageref{31}
\end{frame}

\begin{frame}
    \begin{exampleblock}{問題を解いてみよう}
    \begin{itemize}
        \item 教科書31ページ 問題5-13(1問)
    \end{itemize}
    \end{exampleblock}
\end{frame}

\begin{frame}[fragile]
    \frametitle{信号を二つ以上登録したいとき}
    \begin{lstlisting}[title=２つの信号を登録するときのtemplate.lircd.comf,label=２つの信号を登録するときのtemplate.lircd.comf]
    begin remote
    name fan
    flags RAW_CODES
    eps 30
    aeps 100

    gap 200000
    toggle_bit_mask 0x0

    begin raw_codes

        name onoff
        1207 588 447 838 447 823 475 812 1218…

        <#red#name power
        1364 363 752 743 372 836 164 614 ...#>

    end raw_codes

    end remote
    \end{lstlisting}
    \textpageref{33}
\end{frame}

\begin{frame}[fragile]
    \frametitle{赤外線受信のまとめ}
    \begin{lstlisting}
    mode2 -d /dev/lirc1 | tee <filename1>
    convert_pattern <filename1> > <filename2>
    <filename3>.confを<filename2>を使って作成
    sudo cp <filename3>.conf /etc/lirc/lircd.d/
    sudo service lircd restart
    \end{lstlisting}
    \textpageref{32}
\end{frame}

\begin{frame}
    \frametitle{赤外線を送るコマンド}
    irsend \textless 送り方\textgreater \textless リモコンの名前\textgreater \textless 信号の名前\textgreater

    \begin{itemize}
        \item \textless 送り方\textgreater
        \begin{itemize}
            \item SEND\_ONCE 一度だけ送る
            \item SEND\_START 送り続ける
            \item SEND\_STOP "SEND\_START"で送り続けていた信号を止める
        \end{itemize}
        \item \textless リモコンの名前\textgreater
        \begin{itemize}
            \item なんとか.lircd.confファイルの"begin remote"の直後の"name"
        \end{itemize}
        \item \textless 信号の名前\textgreater
        \begin{itemize}
            \item なんとか.lircd.confファイルの"begin codes"または"begin codes\_raw"の直後の"name"
        \end{itemize}
    \end{itemize}
\end{frame}

\begin{frame}
    \frametitle{センサー紹介}
    教科書34ページから46ページに詳しいセンサーの紹介があります。
\end{frame}

\begin{frame}
    \frametitle{宿題}
    教科書32ページ 問題 5-14
    \begin{itemize}
        \item 自分の家にある家電をRaspberry Piで制御してみましょう。
        \item エアコン、扇風機、テレビなど赤外線を使って操作する家電を選びましょう。
        \item 赤外線で動かない家電もあります。宿題の結果が動かなかった、でもOKです。
        \item 1問
    \end{itemize}
\end{frame}\end{document}
